\section{Method}

\subsection{Muon matrix update}
ORBIT retains Muon as the base update for matrix-valued parameters. For a matrix parameter
$W_t$ with gradient $G_t$ and momentum state $m_t$, the implementation forms
\begin{equation}
m_t=\mu m_{t-1}+G_t,\qquad
U_t=\alpha(W_t)\,\operatorname{NS}_5\!\left(G_t+\mu m_t\right),
\label{eq:muon-candidate}
\end{equation}
where $\mu=0.95$, $\operatorname{NS}_5$ denotes five iterations of the quintic
Newton--Schulz polar map, and
$\alpha(W)=\sqrt{\max(1,n_{\mathrm{row}}/n_{\mathrm{col}})}$ is the standard
aspect-ratio scaling used by the Muon implementation~\cite{jordan2024muon}. ORBIT changes
this candidate only for query and key projections. Other two-dimensional hidden operators
use the ordinary Muon update, while embeddings, the tied language-model head, biases, and
normalization vectors use the auxiliary AdamW path described in
Section~\ref{sec:training-setup}.

\subsection{RoPE-conditioned query--key geometry}
For a single attention head and RoPE frequency pair $f$, the pre-softmax score contribution
between positions $i$ and $j$ is
\begin{equation}
  s_{ij,f}=q_{i,f}^{\top}R_f(\Delta)k_{j,f},\qquad \Delta=i-j.
  \label{eq:rope-score}
\end{equation}
A perturbation of the query therefore interacts with the rotated key statistics, while a
key perturbation interacts with the corresponding rotated query statistics. Let
$C_{Q,f}$ and $C_{K,f}$ denote exponential-moving-average $2\times2$ covariances of the
unrotated query and key coordinates. ORBIT defines
\begin{align}
M_{Q,f} &= \mathbb{E}_{\Delta}\!\left[R_f(\Delta)C_{K,f}R_f(\Delta)^\top\right],\\
M_{K,f} &= \mathbb{E}_{\Delta}\!\left[R_f(\Delta)^\top C_{Q,f}R_f(\Delta)\right].
\label{eq:orbit-metrics}
\end{align}
The expectation is evaluated on the displacement grid
$\mathcal{D}=\{1,2,4,8,16,32,64,128\}$. Covariances are updated during training with EMA
coefficient $\beta=0.95$ after the first observation. These statistics are optimizer-only
buffers and are absent from the inference graph.

\subsection{ORBIT preconditioning}
Let $U_Q$ and $U_K$ be the Muon candidates from Equation~\ref{eq:muon-candidate}. For each
head and RoPE frequency pair,
\begin{align}
\widehat U_{Q,f} &= M_{Q,f}^{-1/2}U_{Q,f},\\
\widehat U_{K,f} &= M_{K,f}^{-1/2}U_{K,f}.
\end{align}
The metrics are regularized by $10^{-5}I$, their effective condition number is capped at
100, and the inverse square root is evaluated analytically for each symmetric positive
$2\times2$ matrix.

A single joint restoration factor then preserves the total Q/K update budget:
\begin{equation}
\rho =
\sqrt{
\frac{\|U_Q\|_F^2+\|U_K\|_F^2}
     {\|\widehat U_Q\|_F^2+\|\widehat U_K\|_F^2}
},
\qquad
\widetilde U_Q=\rho\widehat U_Q,\quad
\widetilde U_K=\rho\widehat U_K.
\label{eq:joint-restore}
\end{equation}
The functional metric therefore changes the direction of the joint Q/K update without
increasing its aggregate Frobenius norm.

Figure~\ref{fig:orbit-overview} summarizes the data flow of the optimizer before the
step-by-step specification in Algorithm~\ref{alg:orbit}.

\IfFileExists{figures/orbit_overview.pdf}{%
\begin{figure}[t]
\centering
\includegraphics[width=0.98\linewidth]{orbit_overview.pdf}
\caption{ORBIT update pathway. Muon candidate directions are reshaped using Q/K covariance
statistics transported through RoPE relative-position rotations, producing frequency-local
$2\times2$ metrics. Analytic inverse-square-root preconditioning is followed by joint
Frobenius-norm restoration before the parameter update. Other parameter classes retain
their standard Muon or auxiliary AdamW path.}
\label{fig:orbit-overview}
\end{figure}}{}

Algorithm~\ref{alg:orbit} gives the corresponding update for one attention layer. Standard
Muon and auxiliary AdamW updates proceed unchanged for the remaining parameter classes.

\begin{algorithm}[t]
\caption{ORBIT update for a RoPE query--key projection pair}
\label{alg:orbit}
\DontPrintSemicolon
\KwIn{$W_Q,W_K$; gradients $G_Q,G_K$; momentum $m_Q,m_K$; covariances $C_Q,C_K$;
learning rate $\eta$; weight decay $\lambda$; momentum $\mu$; EMA $\beta$;
displacements $\mathcal{D}$}
\textbf{Forward statistics:} update $C_Q,C_K$ from current unrotated Q/K activations
with EMA coefficient $\beta$\;
$m_Q \leftarrow \mu m_Q+G_Q$;\quad $m_K \leftarrow \mu m_K+G_K$\;
$U_Q \leftarrow \operatorname{MuonCandidate}(G_Q+\mu m_Q)$;\quad
$U_K \leftarrow \operatorname{MuonCandidate}(G_K+\mu m_K)$\;
\ForEach{attention head $h$ and RoPE frequency pair $f$}{
  $M_{Q,f}\leftarrow |\mathcal{D}|^{-1}\sum_{\Delta\in\mathcal{D}}
  R_f(\Delta)C_{K,f}R_f(\Delta)^\top + 10^{-5}I$\;
  $M_{K,f}\leftarrow |\mathcal{D}|^{-1}\sum_{\Delta\in\mathcal{D}}
  R_f(\Delta)^\top C_{Q,f}R_f(\Delta) + 10^{-5}I$\;
  $T_{Q,f}\leftarrow \operatorname{InvSqrt}_{2\times2}(M_{Q,f})$;\quad
  $T_{K,f}\leftarrow \operatorname{InvSqrt}_{2\times2}(M_{K,f})$\;
  $\widehat U_{Q,f}\leftarrow T_{Q,f}U_{Q,f}$;\quad
  $\widehat U_{K,f}\leftarrow T_{K,f}U_{K,f}$\;
}
$\rho\leftarrow
\sqrt{(\|U_Q\|_F^2+\|U_K\|_F^2)/
(\|\widehat U_Q\|_F^2+\|\widehat U_K\|_F^2)}$\;
$\widetilde U_Q\leftarrow\rho\widehat U_Q$;\quad
$\widetilde U_K\leftarrow\rho\widehat U_K$\;
$W_Q\leftarrow(1-\eta\lambda)W_Q-\eta\widetilde U_Q$;\quad
$W_K\leftarrow(1-\eta\lambda)W_K-\eta\widetilde U_K$\;
\end{algorithm}

\subsection{Cost and ablation variants}
The additional state consists of two $2\times2$ covariance matrices for every attention
head and RoPE frequency pair. The unique per-step matrix operations are local $2\times2$
transforms rather than dense model-width preconditioners, leaving inference unchanged.

Three controls isolate the update components. \textbf{Identity} keeps the ORBIT training
path and statistics collection but disables functional preconditioning. \textbf{No-RoPE}
uses Q/K covariance metrics without relative-position rotation. \textbf{Diagonal} retains
RoPE-aware per-frequency scaling while removing off-diagonal coupling within each
$2\times2$ pair.

\section{Experimental Setup}
\subsection{Models and data}
All reported comparisons use the same decoder implementation within a given size. The
124M model has 12 layers, 12 attention heads, width 768, and approximately 124M parameters.
The 355M model has 24 layers, 16 heads, width 1024, and uses gradient checkpointing. Both
use context length 512, GPT-2 token embeddings with tied input/output weights, pre-norm
Transformer blocks, GELU MLPs with $4d$ hidden width, and zero dropout. Learned absolute
position embeddings are replaced by RoPE with base 10,000. Attention uses separate
$q_{\mathrm{proj}}$, $k_{\mathrm{proj}}$, and $v_{\mathrm{proj}}$ matrices.

The corpus is FineWeb-Edu. Text is tokenized with the GPT-2
tokenizer. A fixed stream of 12.4M tokens is shared across runs: the first 400,000 tokens
form the validation pool and the next 12,000,000 tokens form the training pool. Every training step
samples random contiguous 512-token windows from that pool with a seeded generator;
evaluation uses a fixed, separate generator and averages 20 validation batches. Thus every
optimizer in a paired setting sees the same validation windows.

\subsection{Training setup}
\label{sec:training-setup}
All optimizers use the same warmup--cosine multiplier. The first 10\% of steps are warmup;
the remaining schedule follows a cosine decay ending at $0.1$ times the peak learning rate.
The same multiplier is applied to every learning-rate field, including the auxiliary AdamW
path. Global gradient norm is clipped at 1.0. Training uses CUDA autocast in float16 with
gradient scaling. The 124M model uses batch size 8 (4,096 sampled tokens per optimizer
step); the 355M model uses batch size 4 (2,048 sampled tokens per step). The latter
difference is disclosed because it prevents interpreting the 355M result as a pure
model-size scaling experiment.

\subsection{Baselines and hyperparameter selection}
The broad campaign includes AdamW, Muon, NorMuon, AdaMuon, ASTRO, and ORBIT. ASTRO is a
pre-existing frozen baseline selected during the preceding ASTRO development study before
any ORBIT result was observed; it is not retuned during the ORBIT campaign.

Every optimizer receives the same number of tuned dimensions in the broad campaign.
Muon-family methods search learning rate in $[2\times10^{-3},10^{-1}]$, weight decay in
$[10^{-3},3\times10^{-1}]$, and auxiliary AdamW multiplier in $[0.02,1.5]$, with
log-uniform deterministic sampling. AdamW tunes learning rate, weight decay, and $\beta_2$.
AdaMuon uses its published natural learning-rate range but the same weight-decay and
auxiliary-rate policy. Five discovery trials are used per optimizer in the broad campaign.

The primary Muon--ORBIT comparison uses a stricter matched policy: ten deterministic
900-step candidate configurations are generated once from the Muon search space and
presented identically to both methods. Both methods independently selected the same candidate from the shared grid before any
held-out matched-confirmation seed was evaluated. That selected setting is reused in
the matched confirmation and expanded ORBIT ablation.

\subsection{Evaluation protocol}
The evaluation is organized into complementary study blocks rather than a single benchmark:
broad optimizer context, matched mechanism isolation, component ablations, and transfer
across training horizon and model size. Table~\ref{tab:protocol} summarizes the design.

\begin{table}[t]
\centering
\caption{Experimental design. The primary estimate is the matched Muon--ORBIT comparison;
the remaining experiments test recipe interaction, mechanism attribution, or transfer.}
\label{tab:protocol}
\small
\setlength{\tabcolsep}{4pt}
\renewcommand{\arraystretch}{1.12}
\begin{tabularx}{\textwidth}{@{}
>{\raggedright\arraybackslash}p{2.05cm}
>{\raggedright\arraybackslash}X
>{\centering\arraybackslash}p{0.95cm}
>{\centering\arraybackslash}p{0.85cm}
>{\centering\arraybackslash}p{1.35cm}
>{\raggedright\arraybackslash}p{3.1cm}
@{}}
\toprule
Experiment & Scientific purpose & Model & Steps & Replication & Compared methods \\
\midrule
Broad benchmark &
Frozen-recipe optimizer comparison &
124M & 900 & $n=5$ &
AdamW, AdaMuon, Muon, NorMuon, ASTRO, ORBIT \\

Crossed recipes &
Separate update-rule and recipe effects &
124M & 900 & $n=5$/cell &
Muon, ORBIT \\

\textbf{Matched confirmation} &
\textbf{Mechanism-isolated primary comparison} &
\textbf{124M} & \textbf{900} & \textbf{$n=10$ paired} &
\textbf{Muon, ORBIT} \\

Ablation &
Attribute the functional, RoPE, and pair-coupling components &
124M & 900 & $n=10$/control &
ORBIT, identity, no-RoPE, diagonal \\

Long-horizon transfer &
Test persistence under longer optimization &
124M & 2700 & $n=2$ &
Muon, NorMuon, ASTRO, ORBIT \\

Scale transfer &
Test transfer to a larger decoder &
355M & 900 & $n=2$ &
Muon, NorMuon, ASTRO, ORBIT \\
\bottomrule
\end{tabularx}

\vspace{0.35em}
\footnotesize
Broad hyperparameter selection used five discovery candidates per optimizer. The matched
comparison used ten shared candidate configurations selected before the held-out runs. The
final ablation followed a three-run pilot and used ten paired runs per control.
\end{table}

At 124M, 300, 900, and 2700 optimizer steps correspond to approximately 1.23M, 3.69M,
and 11.06M sampled training tokens for the discovery, primary, and long-horizon regimes.
The 355M transfer uses 1.84M sampled training tokens at 900 steps because of its smaller
batch size.

\paragraph{Statistical analysis.}
The primary estimand is the paired per-run validation-loss difference, ORBIT minus the
comparison method. Reported statistics include the paired mean, sample standard deviation,
paired-run win count, and a two-sided 95\% Student-$t$ confidence interval for the mean
paired difference. The matched Muon--ORBIT comparison is the primary confirmatory contrast;
ablations provide mechanism attribution, while the long-horizon and 355M experiments assess
transfer. Wall-clock time and peak GPU-memory allocation are reported alongside loss.

All final comparisons use the same frozen implementation and a consistent run-record format.
Optimizer settings, random seeds, environment metadata, validation metrics, runtime,
memory use, and training trajectories are retained with each run. Results are merged only
after automated consistency checks verify the expected experiment structure, matched
configuration, finite metrics, and a common implementation/software state. Exact
environment details and run metadata are retained in the accompanying reproducibility
artifacts. The optimizer implementation, frozen run records, and reproduction scripts are
publicly available at \href{https://github.com/pop123-ux/ORBIT}{github.com/pop123-ux/ORBIT}.
